\documentclass[12pt]{article}
\usepackage{amsmath,amssymb,amsfonts}
\usepackage[margin=1in]{geometry}

\begin{document}

\begin{center}
{\Large\bfseries Chapter 10, Problem 22}
\end{center}

\noindent\textbf{Problem.} In quantum mechanics one defines an electron spin wave function using $\alpha$ for spin ``up'' and $\beta$ for spin ``down''; $\alpha$ and $\beta$ are orthonormal. For a system of three electrons we can form product wave functions of the form $\alpha(1)\alpha(2)\alpha(3)$, $\alpha(1)\alpha(2)\beta(3)$, etc., there being $2^3 = 8$ such products. Clearly, these eight functions transform among themselves under the action of the elements of $S_3$, and they are mutually orthogonal. They therefore form a set of basis functions for a representation of $S_3$. Find how many times the various irreducible representations of $S_3$ occur in this representation. Is any representation missing? Why? Can you relate some of the objects you're dealing with to physics?

\bigskip
\noindent\textbf{Solution.}

The group $S_3$ acts on the three-electron system by permuting the electron labels. The 8 basis functions are all products $\sigma_1(1)\sigma_2(2)\sigma_3(3)$ where $\sigma_i \in \{\alpha, \beta\}$.

\medskip
\noindent\textbf{Character table of $S_3$:}
\[
\begin{array}{c|ccc}
S_3 & \{e\} & \{(12),(13),(23)\} & \{(123),(132)\} \\
& 1 & 3 & 2 \\
\hline
\Gamma_1 & 1 & 1 & 1 \\
\Gamma_2 & 1 & -1 & 1 \\
\Gamma_3 & 2 & 0 & -1
\end{array}
\]

\noindent\textbf{Finding the characters of the 8-dimensional representation:}

$\chi(e) = 8$ (identity fixes all 8 basis functions).

For a transposition like $(12)$: $P_{12}\,\sigma_1(1)\sigma_2(2)\sigma_3(3) = \sigma_1(2)\sigma_2(1)\sigma_3(3) = \sigma_2(1)\sigma_1(2)\sigma_3(3)$. This equals the original function only if $\sigma_1 = \sigma_2$. The fixed functions are: $\alpha\alpha\alpha$, $\alpha\alpha\beta$, $\beta\beta\alpha$, $\beta\beta\beta$ (where the first two indices are equal). So $\chi((12)) = 4$.

For a 3-cycle $(123)$: $P_{123}\,\sigma_1(1)\sigma_2(2)\sigma_3(3) = \sigma_1(3)\sigma_2(1)\sigma_3(2) = \sigma_2(1)\sigma_3(2)\sigma_1(3)$. This equals the original if $\sigma_1 = \sigma_2 = \sigma_3$. Only $\alpha\alpha\alpha$ and $\beta\beta\beta$ are fixed. So $\chi((123)) = 2$.

\medskip
\noindent\textbf{Decomposition:}

$a_\nu = \frac{1}{6}\sum_i g_i \chi_i^{(\nu)*}\chi_i$.

\begin{align*}
a_1 &= \frac{1}{6}[1\cdot 1\cdot 8 + 3\cdot 1\cdot 4 + 2\cdot 1\cdot 2] = \frac{1}{6}[8 + 12 + 4] = 4, \\
a_2 &= \frac{1}{6}[1\cdot 1\cdot 8 + 3\cdot(-1)\cdot 4 + 2\cdot 1\cdot 2] = \frac{1}{6}[8 - 12 + 4] = 0, \\
a_3 &= \frac{1}{6}[1\cdot 2\cdot 8 + 3\cdot 0\cdot 4 + 2\cdot(-1)\cdot 2] = \frac{1}{6}[16 + 0 - 4] = 2.
\end{align*}

Check: $4\cdot 1 + 0\cdot 1 + 2\cdot 2 = 8$. \checkmark

\medskip
\noindent\textbf{Result:}
\[
D^{(8)} = 4\Gamma_1 + 0\Gamma_2 + 2\Gamma_3.
\]

The \textbf{antisymmetric representation $\Gamma_2$ is missing.}

\medskip
\noindent\textbf{Why is $\Gamma_2$ absent?}

The representation $\Gamma_2$ corresponds to the fully antisymmetric (alternating) representation, where every transposition gives a factor of $-1$. A function that is completely antisymmetric under exchange of all three electrons would require $f(1,2,3) = -f(2,1,3) = -f(1,3,2)$, etc. For spin-$\frac{1}{2}$ particles with only two possible spin states ($\alpha, \beta$), it is impossible to antisymmetrize three indices when each takes only two values --- by the pigeonhole principle, at least two electrons must have the same spin, making complete antisymmetrization impossible. (This is the same reason a $3\times 3$ determinant with entries from only two values vanishes.)

\medskip
\noindent\textbf{Physical interpretation:}

The four copies of $\Gamma_1$ (symmetric representation) correspond to the total spin $S = 3/2$ multiplet with $S_z = 3/2, 1/2, -1/2, -3/2$ (the quartet state):
\begin{align*}
|S=3/2, S_z = 3/2\rangle &= \alpha\alpha\alpha, \\
|S=3/2, S_z = 1/2\rangle &= \tfrac{1}{\sqrt{3}}(\alpha\alpha\beta + \alpha\beta\alpha + \beta\alpha\alpha), \\
&\text{etc.}
\end{align*}

The two copies of $\Gamma_3$ (the 2D representation) correspond to two $S = 1/2$ doublet states (each with $S_z = \pm 1/2$), which have mixed symmetry. These are physically relevant in atomic and molecular physics where three-electron systems arise (e.g., the lithium atom).

\end{document}
